\section{A positive theta-kernel family}\label{sec:theta}
Define, for \(u\geq0\),
\begin{align}
 \Phi(u)&=\sum_{n\geq1}\pi n^2e^{5u}
       (2\pi n^2e^{4u}-3)\exp(-\pi n^2e^{4u}),\label{eq:theta}\\
 F(t;\lambda,\mu,\nu)&=\int_0^\infty
       \Phi(u)e^{\lambda u^2+\mu u^4+\nu u^6}\cos(2tu)\dd u .
                                                        \label{eq:F}
\end{align}
Every summand of \(\Phi\) is positive on this domain.
Double exponential decay permits parameter differentiation on compact
sets, including complex parameter sets.
Write \(D_j=\partial_t^jF\). Then
\begin{equation}\label{eq:momentdictionary}
 \partial_\lambda D_j=-D_{j+2}/4,\qquad
 \partial_\mu D_j=D_{j+4}/16,\qquad
 \partial_\nu D_j=-D_{j+6}/64 .
\end{equation}

\subsection{The cusp branch and the design problem}
A validated zero of \((D_0,D_1,D_2)\) at \(\nu=0\) has the rounded
location
\[
 (t,\lambda,\mu)\simeq
 (41.4003413586842508894,\,
  -3.64569206160491909881,\,
   8.33512049891860723771).
\]
These displayed coordinates are explanatory approximations; the
companion certificate specifies the actual enclosing box.
At that zero,
\begin{align*}
 3.33563055\cdot10^{-13}&<D_3<3.33563057\cdot10^{-13},\\
 -9.71679534\cdot10^{-13}&<D_4<-9.71679532\cdot10^{-13}.
\end{align*}
In particular the control determinant for the equations \(D_0=D_1=0\)
is \(D_3D_4/64\ne0\), and \(D_3\ne0\) gives the usual local cusp
geometry of the zero set of \(F\). A computer-assisted continuation
defines a single analytic branch
\[
 Q(\nu)=(\tau(\nu),\lambda(\nu),\mu(\nu)),\qquad
 \nu\in I=[-1/64,0].
\]
Shared parameter faces are joined by root containment in the next
uniqueness box, not merely by overlap of the surrounding boxes.
The validated continuation methods are used to establish the particular
branch and bounds here; cusp validation as a general method is established
in the literature \cite{lessard}.

Differentiating the three cusp equations gives the useful identities
\begin{align}
 \mu'&=D_6/(4D_4),&
 \lambda'&=(4D_5\mu'-D_7)/(16D_3),\nonumber\\
 \tau'&=(D_8/64+D_4\lambda'/4-D_6\mu'/16)/D_3 .\label{eq:velocity}
\end{align}
Set
\[
 w_\nu(u)=\Phi(u)e^{\lambda(\nu)u^2+\mu(\nu)u^4+\nu u^6},\qquad
 p_{j,\nu}(u)=(2u)^j\cos(2\tau(\nu)u+j\pi/2),
\]
and let \(q_{j,\nu}=w_\nu p_{j,\nu}\), \(0\leq j\leq3\).
In \eqref{eq:problem} take \(m=3\) and
\(f_\nu=D_3(Q(\nu),\nu)>0\), and denote its value by \(\delta_\nu(M)\).
All coefficients in Section~\ref{sec:general} now depend on \(\nu\).

If \(g\) is feasible, multiplying the kernel by \(1-g\) makes the
first four \(t\)-derivatives, of orders \(0,1,2,3\), vanish at the
prescribed point. If \(\norm{g}_\infty<1\), the modified kernel is still
positive. The size \(\delta_\nu(M)\) is therefore the least amplitude of
such a constrained change, with its roughness controlled by \(M\).
This interpretation concerns a separate profile for each \(\nu\).

\begin{theorem}[Theta value enclosure]\label{thm:theta}
For every \(\nu\in I\), the sign certificate and hypotheses of
Theorem~\ref{thm:infinite} hold for the kernels just defined.
For every \(M\geq M_0=2\cdot10^{-5}\), the optimum is attained and
\[
 9.17\cdot10^{-10}<\delta_\nu(M)<9.181\cdot10^{-10}<1 .
\]
With \(P_{4,\nu}(M)=\delta_{0,\nu}+C_{2,\nu}M^{-2}+C_{4,\nu}M^{-4}\),
\begin{equation}\label{eq:thetaerror}
 -\frac{1.02\cdot10^{-53}}{M^6}
 <\delta_\nu(M)-P_{4,\nu}(M)
 <\frac{1.64\cdot10^{-53}}{M^6}.
\end{equation}
Furthermore, throughout \(I\),
\begin{align}
 2.48374879\cdot10^{-39}&<C_{4,\nu}<2.49203006\cdot10^{-39},
                                                        \label{eq:C4range}\\
 2.7\cdot10^{-11}&<\delta_{0,\nu}'<3.0\cdot10^{-11},\nonumber\\
 7.4\cdot10^{-26}&<C_{2,\nu}'<8.7\cdot10^{-26},\nonumber\\
 1.9\cdot10^{-40}&<C_{4,\nu}'<5.3\cdot10^{-40}.\label{eq:coefficientmotion}
\end{align}
\end{theorem}

\begin{proof}[Proof structure and computational evidence]
The branch construction, sign equations, invertible moment Jacobian,
simple residual zeros, and derivative bounds are established on
32 closed driver cells
\[
 I_i=[-(i+1)/2048,-i/2048],\qquad 0\leq i<32.
\]
The moment Jacobian is \(-G\); its nonsingularity and the convexity of
\(\int|q_b|\) identify the local dual solution as the unique global
minimizer. For clarity, a second minimizer would make the objective
constant on their joining segment, contradicting its positive
definite Hessian near the certified minimizer.

The residual without its positive weight has the form
\begin{equation}\label{eq:rawresidual}
 r_b(u)=(8u^3+2b_1u)\sin(2\tau u)
                  +(4b_2u^2-b_0)\cos(2\tau u).
\end{equation}
On \([0,1]\), the evidence isolates 28 simple zeros and covers their
complement by sign enclosures. On \(u\geq1\), put
\(T=( (8u^3+2b_1u)^2+(4b_2u^2-b_0)^2)^{1/2}\) and
\(\phi=2\tau u+\arctan((4b_2u^2-b_0)/(8u^3+2b_1u))\).
The uniform bounds
\[
 T>7u^3,\quad |T'|<25u^2,\quad 81<\phi'<85
\]
give separated, simple zeros, with spacing greater than \(3/85\) and
\(|r_b'|>567u^3\) at each zero.
They also bound root motion with \(b\), since
\(|\partial_{b_j}z|\leq1/(81z)\) for the tail zeros.

The weight decay verifies the stronger, weighted hypothesis
\eqref{eq:weighted}, rather than only unweighted convergence.
For \(H=1/64\) and \(\rho\leq1/1000\), suitable uniform constants give
\begin{align*}
 E_k&\leq C(1+z_k+\rho)^{18}
 e^{21(z_k+\rho)+9(z_k+\rho)^4-\pi e^{4(z_k-\rho)}},\\
 L_k&=e^{-4(z_k+\rho)^2-H(z_k+\rho)^6-\pi e^{4(z_k+\rho)}}.
\end{align*}
The ratio \(E_k^3/L_k^2\) is bounded by a polynomial times
\[
 \exp\{63(z_k+\rho)+27(z_k+\rho)^4+8(z_k+\rho)^2
       +2H(z_k+\rho)^6-\pi c_\rho e^{4z_k}\},
 \quad c_\rho=3e^{-4\rho}-2e^{4\rho}>0.97 .
\]
The separated zeros and this double exponential decay prove
\eqref{eq:weighted}. Finite initial neighborhoods are covered directly.
Theorem~\ref{thm:infinite} now applies.

The explicit remainder requires additional estimates and does not
follow just from that theorem. Appendix~\ref{app:effective} constructs
finite ramp profiles with exact moments and an accompanying support
bound, uniformly on all 32 cells. Scalar comparison, without continuity
of the number of active ramps, proves \eqref{eq:thetaerror} and
feasibility for every stated \(M\).
Differentiated coefficient identities and interval enclosures give
\eqref{eq:C4range}--\eqref{eq:coefficientmotion}.
The exact input-output map and replay commands appear in
Appendix~\ref{app:evidence}; those finite computations are part of this
computer-assisted theorem.
\end{proof}

At \(M=M_0\), the absolute error in \eqref{eq:thetaerror} lies between
\(-1.59375\cdot10^{-25}\) and \(2.5625\cdot10^{-25}\).
These figures measure the truncation error once the coefficients are
known; rounding the coefficients separately introduces additional
error. In particular the broad uniform ranges above are not substitutes
for pointwise coefficient enclosures.
